\section{Disseminazione}

\textbf{PhD Workshop UnaEuropa}: ``Museums and New Challenges: virtual technologies, social responsibility and environmental sustainability''\\
\textit{Luogo}: Bologna, Italia\\
\textit{Data}: 5-10 maggio 2025\\
\textit{Descrizione}: Workshop internazionale per dottorandi nell'ambito del patrimonio culturale, focalizzato sulle nuove sfide che i musei devono affrontare in termini di tecnologie virtuali, responsabilità sociale e sostenibilità ambientale.\\
\textit{Presentazione}: ``HERITRACE: A User-Friendly Semantic Data Editor with Change Tracking and Provenance Management for Cultural Heritage Institutions''. Presentazione di HERITRACE, uno strumento web-based progettato per colmare il divario tra tecnologie semantiche complesse e l'expertise degli operatori del patrimonio culturale. Il sistema presenta un'interfaccia user-friendly per la gestione di dati RDF, gestione completa della provenienza, change-tracking con funzionalità snapshot e personalizzazione flessibile attraverso linguaggi standard come SHACL e YAML.\\
\textit{DOI}: \url{https://doi.org/10.5281/zenodo.15375769}

\vspace{0.5cm}

\textbf{csv,conf,v9}\\
\textit{Luogo}: Bologna, Italia\\
\textit{Data}: 10-11 settembre 2025\\
\textit{Descrizione}: Conferenza internazionale dedicata alla gestione e analisi di dati strutturati.\\
\textit{Presentazione}: ``Mapping the unmapped citation landscape: how crowdsourcing will fill the citation gap''. La presentazione esamina il problema dell'invisibilità delle citazioni nell'editoria accademica, particolarmente per le riviste del Sud Globale. Basandosi su ricerche esistenti che mostrano come tra oltre 52.000 riviste Open Journal Systems (OJS) mondiali, il 79,9\% si trova nel Sud Globale, con il 98,8\% invisibile a Web of Science e il 94,3\% invisibile a Scopus, la presentazione introduce una soluzione collaborativa sviluppata attraverso una partnership tra OpenCitations, Public Knowledge Project (PKP) e Leibniz Information Centre. La soluzione implementa un flusso di lavoro integrato in OJS che consente l'invio automatico di metadati bibliografici e citazioni a OpenCitations.\\
\textit{DOI}: \url{https://doi.org/10.5281/zenodo.17098599}